%───────────────────────────────────────────────────────────────────────────────
% NOTATION.TEX — Vocabulario matemático del documento
%───────────────────────────────────────────────────────────────────────────────
%
% Este archivo define los objetos matemáticos propios del documento.
% NO configura paquetes — eso vive en config/equations.tex y demás.
%
% RELACIÓN CON OTROS ARCHIVOS:
%   physics       →  \abs, \norm, \tr, \grad, \div, \curl, \laplacian,
%                    \dv, \pdv, \bra, \ket, \braket, \comm, \acomm, etc.
%   equations.tex →  \eqop, \equ, \eqk, \eqfn, \eqdm, \eqdom (coloreado semántico)
%                    \eqpdv, \eqgrad, \eqdvg, \eqSDE  (compuestos coloreados)
%   notation.tex  →  todo lo demás: objetos del dominio, operadores propios,
%                    distribuciones, shortcuts de uso frecuente.
%
% REGLA DE EXTENSIÓN:
%   Añadir secciones por dominio. Nunca mezclar dominios dentro de una sección.
%   Si una sección supera ~30 comandos, considerar un archivo propio.
%
% CONFLICTOS CONOCIDOS CON physics:
%   \div   →  physics lo redefine como divergencia. Usar \eqdvg (coloreado)
%             o \divergence (physics). Para división: / o \frac{}{}.
%   \rank  →  physics define \rank. Para sobrescribir: ver nota al final.
%───────────────────────────────────────────────────────────────────────────────


% =============================================================================
% ÁLGEBRA Y ANÁLISIS
% =============================================================================

% ─── Conjuntos numéricos ─────────────────────────────────────────────────────
\newcommand{\R}{\mathbb{R}}
\newcommand{\N}{\mathbb{N}}
\newcommand{\Z}{\mathbb{Z}}
\newcommand{\C}{\mathbb{C}}
\newcommand{\Q}{\mathbb{Q}}
\newcommand{\F}{\mathbb{F}}          % cuerpo genérico (álgebra abstracta)

% ─── Operadores matriciales y lineales ───────────────────────────────────────
% \tr, \rank ya existen en physics. Ver nota al final para sobrescribir \rank.
\newcommand{\T}{^{\top}}             % transpuesta: A\T
\newcommand{\HH}{^{\dagger}}         % adjunto hermítico (†): A\HH
\newcommand{\Id}{\mathbf{I}}         % matriz identidad
\DeclareMathOperator{\diag}{diag}    % \diag(a,b,c) o \diag(A)
\DeclareMathOperator{\spec}{spec}    % espectro de un operador

% ─── Optimización ────────────────────────────────────────────────────────────
\DeclareMathOperator*{\argmin}{arg\,min}
\DeclareMathOperator*{\argmax}{arg\,max}
\DeclareMathOperator*{\minimize}{minimize}
\DeclareMathOperator*{\maximize}{maximize}

% ─── Aliases de physics (nombres cortos para uso frecuente) ──────────────────
% Mantener solo los que se usan en el documento con alta frecuencia.
% Para el resto: usar directamente \laplacian, \divergence, etc.
\newcommand{\Lap}{\laplacian}        % alias → \laplacian (∇²)
\newcommand{\Div}{\divergence}       % alias → \divergence (∇·)


% =============================================================================
% PROBABILIDAD Y ESTADÍSTICA
% =============================================================================

% ─── Operadores de esperanza y probabilidad ───────────────────────────────────
% \mleft[\mright] corrige el espaciado que \left[\right] introduce en
% ciertos contextos (ver paquete mleftright).
\newcommand{\E}[1]{\mathbb{E}\!\mleft[#1\mright]}
\newcommand{\Prob}[1]{\mathbb{P}\!\mleft(#1\mright)}
\newcommand{\Ep}[1]{\langle #1 \rangle}   % valor esperado estilo física ⟨·⟩

% ─── Operadores estadísticos ─────────────────────────────────────────────────
\DeclareMathOperator{\KL}{KL}
\DeclareMathOperator{\Var}{Var}
\DeclareMathOperator{\Cov}{Cov}
\DeclareMathOperator{\Ent}{H}        % entropía: H(X)
\DeclareMathOperator{\MI}{I}         % información mutua: I(X;Y)

% ─── Distribuciones ──────────────────────────────────────────────────────────
\newcommand{\Normal}{\mathcal{N}}    % \Normal(\mu, \sigma^2)
\newcommand{\Uniform}{\mathcal{U}}   % \Uniform(a, b)
\newcommand{\Bernoulli}{\mathrm{Ber}}
\newcommand{\Poisson}{\mathrm{Poi}}
\newcommand{\Categorical}{\mathrm{Cat}}


% =============================================================================
% MACHINE LEARNING
% =============================================================================

% ─── Objetos del modelo ──────────────────────────────────────────────────────
\newcommand{\Loss}{\mathcal{L}}      % función de pérdida ℒ
\newcommand{\Lscr}{\mathscr{L}}      % pérdida en notación script (alternativa)
\newcommand{\D}{\mathcal{D}}         % dataset 𝒟
\newcommand{\Params}{\Theta}         % espacio de parámetros Θ
\newcommand{\params}{\theta}         % parámetros θ

% ─── Funciones de activación ─────────────────────────────────────────────────
\DeclareMathOperator{\softmax}{softmax}
\DeclareMathOperator{\sigmoid}{\sigma}   % \sigmoid(x) = σ(x)
\DeclareMathOperator{\relu}{ReLU}
\DeclareMathOperator{\gelu}{GELU}

% ─── Score function y difusión ───────────────────────────────────────────────
% \score toma la variable como argumento opcional (default: \bm{x})
% Uso: \score  →  ∇_x log p(x)
%      \score[\bm{z}]  →  ∇_z log p(z)
\newcommand{\score}[1][\bm{x}]{\grad_{#1}\log p}

% Ecuación diferencial estocástica — versión sin color (inline, B/N)
% Para la versión con coloreado semántico usar \eqSDE de equations.tex.
% \SDE{X}{a(X,t)}{b(t)}  →  dX = a(X,t) dt + b(t) dW
\newcommand{\SDE}[3]{d#1 = #2\,dt + #3\,dW}


% =============================================================================
% FÍSICA COMPUTACIONAL
% =============================================================================

% ─── Mecánica cuántica ───────────────────────────────────────────────────────
% \bra, \ket, \braket, \comm, \acomm ya existen en physics.
\newcommand{\Ham}{\hat{H}}           % operador Hamiltoniano Ĥ
\newcommand{\Hscr}{\mathscr{H}}      % espacio de Hilbert ℋ
\newcommand{\Fscr}{\mathscr{F}}      % transformada de Fourier ℱ (operador)

% ─── Mecánica estadística ─────────────────────────────────────────────────────
\newcommand{\kB}{k_{\mathrm{B}}}     % constante de Boltzmann k_B
\newcommand{\Zpart}{Z}               % función de partición Z
\newcommand{\FP}{\mathcal{L}_{\mathrm{FP}}}  % operador Fokker-Planck ℒ_FP

% ─── Física de reactores ─────────────────────────────────────────────────────
\newcommand{\keff}{k_{\mathrm{eff}}} % factor de multiplicación efectivo


% =============================================================================
% SÍMBOLOS AUXILIARES
% =============================================================================
\newcommand{\OK}{\ensuremath{\checkmark}}   % ✓
\newcommand{\nOK}{\ensuremath{\times}}      % ✗


% =============================================================================
% NOTAS DE COMPATIBILIDAD
% =============================================================================
%
% ── Sobrescribir \rank de physics ────────────────────────────────────────────
% physics define \rank como operador. Para usar "rango" en español
% o cambiar el comportamiento, descomentar:
% \let\rank\relax
% \DeclareMathOperator{\rank}{rango}
%
% ── \D puede colisionar ──────────────────────────────────────────────────────
% Si algún paquete cargado define \D (poco común), usar \Dataset en su lugar:
% \let\D\relax
% \newcommand{\Dataset}{\mathcal{D}}
%
% ── \sigmoid vs variable σ ───────────────────────────────────────────────────
% \sigmoid usa \sigma internamente. Si σ ya está ocupada como variable en
% una sección, usar \sigma(x) directamente en lugar de \sigmoid(x).
